\chapter{Hilbert Spaces}

Throughout this chapter, the inner product is linear in its first argument and
conjugate linear in its second.

\section{Inner products and their geometry}

\begin{notedefinition}[Inner product]
An inner product on a real or complex vector space $H$ is a map
$\langle\cdot,\cdot\rangle\colon H\times H\to\mathbb K$ satisfying
\begin{enumerate}
  \item $\langle x,x\rangle\ge0$, with equality only for $x=0$;
  \item $\langle \alpha x+\beta y,z\rangle
        =\alpha\langle x,z\rangle+\beta\langle y,z\rangle$;
  \item $\langle x,y\rangle=\overline{\langle y,x\rangle}$.
\end{enumerate}
The associated norm is
\[
  \|x\|=\sqrt{\langle x,x\rangle}.
\]
A complete inner-product space is a Hilbert space.
\end{notedefinition}

\begin{notetheorem}[Cauchy--Schwarz inequality]
For all $x,y\in H$,
\[
  |\langle x,y\rangle|\le\|x\|\,\|y\|.
\]
Equality holds exactly when $x$ and $y$ are linearly dependent.
\end{notetheorem}

\begin{proof}
If $y\ne0$, positivity of
$\|x-\langle x,y\rangle y/\|y\|^2\|^2$ gives the inequality.  Equality is
equivalent to that orthogonal remainder being zero.
\end{proof}

\begin{notedefinition}[Orthogonality]
Vectors $x,y$ are orthogonal, written $x\perp y$, if
$\langle x,y\rangle=0$.
\end{notedefinition}

\begin{notetheorem}[Pythagorean theorem and orthogonal decomposition]
If $x\perp y$, then
\[
  \|x+y\|^2=\|x\|^2+\|y\|^2.
\]
For $y\ne0$ every $x$ has the decomposition
\[
  x=\frac{\langle x,y\rangle}{\|y\|^2}y
    +\left(x-\frac{\langle x,y\rangle}{\|y\|^2}y\right),
\]
where the second term is orthogonal to $y$.
\end{notetheorem}

\begin{notetheorem}[Triangle inequality]
The inner-product norm satisfies
\[
  \|x+y\|\le\|x\|+\|y\|.
\]
Equality holds exactly when one vector is a nonnegative real multiple of the
other.
\end{notetheorem}

\begin{notetheorem}[Parallelogram identity]
For all $x,y\in H$,
\[
  \|x+y\|^2+\|x-y\|^2=2\|x\|^2+2\|y\|^2.
\]
\end{notetheorem}

\begin{noteexample}[$\ell^p$ is not Hilbert when $p\ne2$]
For $x=(1,1,0,\ldots)$ and $y=(1,-1,0,\ldots)$,
\[
  \|x+y\|_p^2+\|x-y\|_p^2=8,
  \qquad
  2\|x\|_p^2+2\|y\|_p^2=4\cdot2^{2/p}.
\]
The parallelogram identity holds only for $p=2$, so the $\ell^p$ norm is not
induced by an inner product when $p\ne2$.
\end{noteexample}

\begin{notetheorem}[Polarization identities]
A norm induced by an inner product determines that inner product.  Over
$\mathbb R$,
\[
  \langle x,y\rangle
  =\frac14\bigl(\|x+y\|^2-\|x-y\|^2\bigr).
\]
Over $\mathbb C$ with linearity in the first variable,
\[
  \langle x,y\rangle
  =\frac14\sum_{k=0}^{3}i^k\|x+i^ky\|^2.
\]
\end{notetheorem}

\begin{notelemma}[Continuity of the inner product]
If $x_n\to x$ and $y_n\to y$, then
$\langle x_n,y_n\rangle\to\langle x,y\rangle$.
\end{notelemma}

\begin{proof}
By Cauchy--Schwarz,
\[
 |\langle x_n,y_n\rangle-\langle x,y\rangle|
 \le\|x_n-x\|\|y_n\|+\|x\|\|y_n-y\|.
\]
The sequence $(y_n)$ is bounded.
\end{proof}

\begin{notetheorem}[Operator norm through the inner product]
Let $A\in\Bop{H_1}{H_2}$ between Hilbert spaces.  Then
\[
  \|A\|
  =\sup_{x\ne0,\,y\ne0}
    \frac{|\langle Ax,y\rangle|}{\|x\|\|y\|}.
\]
\end{notetheorem}

\begin{proof}
Cauchy--Schwarz gives the upper bound.  For $x$ near norm attaining and
$y=Ax$, the quotient equals $\|Ax\|/\|x\|$.
\end{proof}

\begin{notetheorem}[Isometries fixing the origin]
Let $H_1,H_2$ be real inner-product spaces and let
$A\colon H_1\to H_2$ satisfy $A0=0$ and
$\|Ax-Ay\|=\|x-y\|$.  Then $A$ is linear and preserves inner products.
\end{notetheorem}

\begin{proof}
Taking $y=0$ gives $\|Ax\|=\|x\|$.  The real polarization identity gives
$\langle Ax,Ay\rangle=\langle x,y\rangle$.  Hence
\[
 \|A(x+y)-Ax-Ay\|^2=0,
\]
so $A$ is additive.  Integer and rational homogeneity follow, and continuity
extends homogeneity to all real scalars.
\end{proof}

\begin{noteexample}[Complex caveat]
Complex conjugation on $\mathbb C$ is an isometry fixing $0$ but is not
complex linear.  Thus the preceding statement is a real theorem unless an
additional complex-linearity condition is imposed.
\end{noteexample}

\begin{notetheorem}[Mazur--Ulam theorem]
Every surjective isometry between real normed spaces is affine.  Thus, if
$f(0)=0$, it is real linear.
\end{notetheorem}

\begin{proof}[Proof sketch]
A reflection argument shows that a surjective isometry preserves midpoints:
$f((x+y)/2)=(f(x)+f(y))/2$.  Iterating yields preservation of dyadic convex
combinations, continuity gives all real convex combinations, and translation
by $f(0)$ turns the map into an additive homogeneous map.  The source's
strictly-convex ball-intersection proof covers only a special case; the
statement above is the full theorem.
\end{proof}

\begin{noteexercise}[Vanishing quadratic form]
Let $H$ be a complex inner-product space and let $T\in\Bop{H}{H}$.  If
\[
  \langle Tx,x\rangle=0\qquad(x\in H),
\]
then $T=0$.  This is false over $\mathbb R$.
\end{noteexercise}

\begin{proof}
Apply complex polarization to the sesquilinear form
$b(x,y)=\langle Tx,y\rangle$.  Its diagonal values determine the entire form,
so $b=0$ and $T=0$.  Over $\mathbb R$, a nonzero skew-symmetric rotation of
$\mathbb R^2$ satisfies $\langle Tx,x\rangle=0$.
\end{proof}

\section{Metric and orthogonal projections}

\begin{notedefinition}[Distance from a set]
For a nonempty subset $C$ of a normed space,
\[
  \dist(x,C)=\inf_{c\in C}\|x-c\|.
\]
A set is convex if it contains every line segment joining two of its points.
\end{notedefinition}

\begin{notetheorem}[Projection theorem for closed convex sets]
Let $C$ be a nonempty closed convex subset of a Hilbert space $H$.  For every
$x\in H$ there is a unique $P_Cx\in C$ such that
\[
  \|x-P_Cx\|=\dist(x,C).
\]
\end{notetheorem}

\begin{proof}
Let $d=\dist(x,C)$ and choose $c_n\in C$ with $\|x-c_n\|\to d$.  The
parallelogram identity gives
\[
 \|c_n-c_m\|^2
 =2\|x-c_n\|^2+2\|x-c_m\|^2
  -4\left\|x-\frac{c_n+c_m}{2}\right\|^2.
\]
The midpoint belongs to $C$, so the last norm is at least $d$; hence $(c_n)$
is Cauchy.  Completeness and closedness give a minimizer.  If $c,c'$ are two
minimizers, the same identity applied to them gives $\|c-c'\|=0$.
\end{proof}

\begin{notedefinition}[Orthogonal projection]
If $M$ is a closed subspace of a Hilbert space, the nearest-point map
$P_M\colon H\to M$ is the orthogonal projection onto $M$.
\end{notedefinition}

\begin{notetheorem}[Characterization of orthogonal projection]
Let $M$ be a closed subspace and $x\in H$.  For $m\in M$, the following are
equivalent:
\begin{enumerate}
  \item $m=P_Mx$;
  \item $x-m\perp M$.
\end{enumerate}
Moreover, $P_M$ is linear,
\[
  \|P_Mx\|\le\|x\|,
\]
and equality holds if and only if $x\in M$.
\end{notetheorem}

\begin{proof}
If $m$ is minimizing, then for every $y\in M$ and scalar $t$,
\[
  \|x-m\|^2\le\|x-m-ty\|^2.
\]
Varying real and imaginary $t$ forces $\langle x-m,y\rangle=0$.
Conversely, Pythagoras gives
$\|x-y\|^2=\|x-m\|^2+\|m-y\|^2$ for $y\in M$.
The orthogonality characterization proves linearity.  Finally,
$\|x\|^2=\|P_Mx\|^2+\|x-P_Mx\|^2$.
\end{proof}

\begin{notetheorem}[Variational inequality for convex projection]
Let $C$ be nonempty, closed, and convex.  For $x\in H$ and $y\in C$,
\[
  y=P_Cx
  \quad\Longleftrightarrow\quad
  \operatorname{Re}\langle x-y,z-y\rangle\le0
  \quad(z\in C).
\]
The projection is nonexpansive:
\[
  \|P_Cx-P_Cu\|\le\|x-u\|.
\]
It is linear if and only if $C$ is a closed subspace.
\end{notetheorem}

\begin{proof}
For $z\in C$, minimality against $y+t(z-y)$ and division by $t>0$ gives the
variational inequality as $t\downarrow0$.  Conversely, expand
\[
 \|x-z\|^2=\|x-y\|^2+\|z-y\|^2
            -2\operatorname{Re}\langle x-y,z-y\rangle.
\]
For $y=P_Cx$ and $v=P_Cu$, apply the inequality twice to obtain
\[
  \|y-v\|^2\le\operatorname{Re}\langle x-u,y-v\rangle
  \le\|x-u\|\|y-v\|.
\]
If $C$ is a subspace, this is the orthogonal projection and is linear.  The
range of any linear projection is a subspace, giving the converse.
\end{proof}

\begin{noteexercise}[Resolvent of a projection]
Let $P$ be a bounded projection, $P^2=P$.  If $\lambda\notin\{0,1\}$, then
$\lambda I-P$ is invertible and
\[
  (\lambda I-P)^{-1}
  =\frac1\lambda(I-P)+\frac1{\lambda-1}P.
\]
\end{noteexercise}

\begin{proof}
Decompose $x=(I-P)x+Px$ into $\Ker P\oplus\Ran P$.  On these two subspaces,
$\lambda I-P$ acts as multiplication by $\lambda$ and $\lambda-1$,
respectively.
\end{proof}

\begin{noteexercise}[Nested convex projections]
Let $C_1\supseteq C_2\supseteq\cdots$ be nonempty closed convex subsets of a
Hilbert space, with nonempty intersection $C$.  If $x_n=P_{C_n}x$, then
$x_n\to P_Cx$.
\end{noteexercise}

\begin{proof}
The distances $\|x-x_n\|$ increase and are bounded above by
$\dist(x,C)$.  The projection inequalities show $(x_n)$ is Cauchy.  Its limit
lies in every $C_n$ and attains the distance to $C$.
\end{proof}

\begin{noteremark}
Source pages 136--137 are blank and contain no mathematical material.
\end{noteremark}

\section{Orthogonal complements and Riesz representation}

\begin{notedefinition}[Orthogonal complement]
For $M\subseteq H$, define
\[
  M^\perp=\{x\in H:\langle x,m\rangle=0\text{ for every }m\in M\}.
\]
\end{notedefinition}

\begin{notetheorem}[Basic properties]
For subsets and subspaces of a Hilbert space:
\begin{enumerate}
  \item $M^\perp$ is a closed subspace;
  \item $M\cap M^\perp=\{0\}$ when $M$ is a subspace;
  \item $M\subseteq N$ implies $N^\perp\subseteq M^\perp$;
  \item $(\overline M)^\perp=M^\perp$;
  \item $M\subseteq(M^\perp)^\perp$.
\end{enumerate}
\end{notetheorem}

\begin{proof}
Linearity is immediate.  Closedness follows by passing inner products to the
limit.  The remaining statements follow directly from the definition.
\end{proof}

\begin{notetheorem}[Double orthogonal complement]
For every subspace $M$ of a Hilbert space,
\[
  (M^\perp)^\perp=\overline M.
\]
Thus $M$ is dense if and only if $M^\perp=\{0\}$.
\end{notetheorem}

\begin{proof}
The inclusion $\overline M\subseteq(M^\perp)^\perp$ is automatic.  If
$x\in(M^\perp)^\perp$, project $x$ onto $\overline M$.  The remainder
$x-P_{\overline M}x$ belongs both to $M^\perp$ and to its orthogonal
complement, hence is zero.
\end{proof}

\begin{notetheorem}[Orthogonal decomposition]
If $M$ is a closed subspace of a Hilbert space $H$, then
\[
  H=M\oplus M^\perp.
\]
Every $x$ has the unique form
\[
  x=P_Mx+(I-P_M)x,
\]
with the two summands in $M$ and $M^\perp$.
\end{notetheorem}

\begin{notetheorem}[Range and kernel of an orthogonal projection]
For a closed subspace $M$,
\[
  \Ran P_M=M,
  \qquad
  \Ker P_M=M^\perp,
  \qquad
  P_{M^\perp}=I-P_M.
\]
\end{notetheorem}

\begin{notetheorem}[Riesz representation theorem]
For every bounded linear functional $\varphi$ on a Hilbert space $H$, there
is a unique $h\in H$ such that
\[
  \varphi(x)=\langle x,h\rangle
  \qquad(x\in H).
\]
Moreover,
\[
  \|\varphi\|=\|h\|.
\]
\end{notetheorem}

\begin{proof}
If $\varphi=0$, take $h=0$.  Otherwise, $N=\Ker\varphi$ is a proper closed
subspace.  Choose a unit vector $u\in N^\perp$ and set
$h=\overline{\varphi(u)}u$; this conjugate is required by the convention that
the inner product is linear in the first variable.  Every $x$ decomposes as
$n+\alpha u$, and then
\[
  \langle x,h\rangle
  =\alpha\varphi(u)=\varphi(x).
\]
Uniqueness follows by testing the difference against itself.  Cauchy--Schwarz
gives $\|\varphi\|\le\|h\|$, while testing $h/\|h\|$ gives equality.
\end{proof}

\begin{notetheorem}[Representation of bounded sesquilinear forms]
Let $H_1,H_2$ be Hilbert spaces and let
$b\colon H_1\times H_2\to\mathbb K$ be bounded, linear in the first variable
and conjugate linear in the second.  Then there is a unique
$S\in\Bop{H_1}{H_2}$ such that
\[
  b(x,y)=\langle Sx,y\rangle,
  \qquad
  \|S\|=\|b\|.
\]
\end{notetheorem}

\begin{proof}
For each $x$, apply Riesz representation to the linear functional
$y\mapsto\overline{b(x,y)}$.  It gives a unique $Sx\in H_2$ such that
$\overline{b(x,y)}=\langle y,Sx\rangle$, equivalently
$b(x,y)=\langle Sx,y\rangle$.  Uniqueness gives linearity in $x$, and the
norm estimate gives boundedness and equality of norms.
\end{proof}

\begin{notecorollary}[Hilbert direct sum]
For every closed subspace $M$,
\[
  H=M\oplus M^\perp.
\]
\end{notecorollary}

\begin{noteexercise}[Strict convexity of the Hilbert ball]
The closed unit ball of an inner-product space is strictly convex: if
$x,y\in S_H$ and $x\ne y$, then
\[
  \left\|\frac{x+y}{2}\right\|<1.
\]
This can fail in a general normed space; for example in $\ell^\infty_2$ the
unit sphere contains line segments.
\end{noteexercise}

\begin{proof}
The parallelogram identity gives
\[
  \left\|\frac{x+y}{2}\right\|^2
  =1-\frac14\|x-y\|^2.
\]
\end{proof}

\begin{notetheorem}[Hilbert-space extension of operators]
Let $U$ be a closed subspace of a Hilbert space $H$, let $W$ be Banach, and
let $S\in\Bop{U}{W}$.  Then
\[
  T=S\circ P_U
\]
is a bounded extension of $S$ to $H$ and
\[
  \|T\|=\|S\|.
\]
For a nonclosed $U$, first extend $S$ continuously to $\overline U$.
\end{notetheorem}

\begin{proof}
Since $P_U|_U=I$ and $\|P_U\|\le1$, one has
$\|T\|\le\|S\|$.  Restriction to $U$ gives the reverse inequality.
\end{proof}

\begin{noteexercise}[Equality in the triangle inequality for functionals]
Let $\varphi,\psi\in H^*$ be nonzero and suppose
\[
  \|\varphi+\psi\|=\|\varphi\|+\|\psi\|.
\]
Then one is a nonnegative real multiple of the other.  The conclusion need not
hold in an arbitrary Banach space.
\end{noteexercise}

\begin{proof}
Represent the functionals by $h,k\in H$.  The hypothesis becomes equality in
the Hilbert-space triangle inequality for $h+k$, so $h$ and $k$ are
nonnegative real multiples.  In $\ell^\infty_2{}^*=\ell^1_2$, equality can
hold for independent positive coordinate functionals.
\end{proof}

\section{Orthonormal families and bases}

\begin{notedefinition}[Orthonormal family]
A family $(e_i)_{i\in I}$ is orthonormal if
\[
  \langle e_i,e_j\rangle=\delta_{ij}.
\]
For finite scalar families,
\[
  \left\|\sum_{i\in F}a_ie_i\right\|^2
  =\sum_{i\in F}|a_i|^2.
\]
\end{notedefinition}

\begin{notedefinition}[Unordered sum]
For a family $(x_i)_{i\in I}$ in a normed space, the unordered sum
$\sum_{i\in I}x_i$ converges to $x$ if for every $\varepsilon>0$ there is a
finite $F_0\subseteq I$ such that
\[
  \left\|x-\sum_{i\in F}x_i\right\|<\varepsilon
\]
for every finite $F\supseteq F_0$.
\end{notedefinition}

\begin{notetheorem}[Linear combinations of an orthonormal family]
Let $(e_i)_{i\in I}$ be orthonormal in a Hilbert space.  The unordered series
\[
  \sum_{i\in I}a_ie_i
\]
converges if and only if
\[
  \sum_{i\in I}|a_i|^2<\infty.
\]
When it converges,
\[
  \left\|\sum_{i\in I}a_ie_i\right\|^2
  =\sum_{i\in I}|a_i|^2.
\]
\end{notetheorem}

\begin{proof}
Differences of finite partial sums have squared norm equal to the corresponding
finite sum of $|a_i|^2$.  Thus the net of finite partial sums is Cauchy exactly
when the scalar unordered sum is finite.  Completeness gives convergence and
passing to the limit gives the norm identity.  Only countably many $a_i$ can
be nonzero.
\end{proof}

\begin{notetheorem}[Bessel's inequality]
For an orthonormal family $(e_i)_{i\in I}$ and $x\in H$,
\[
  \sum_{i\in I}|\langle x,e_i\rangle|^2\le\|x\|^2.
\]
\end{notetheorem}

\begin{proof}
For a finite $F$, let
$y_F=\sum_{i\in F}\langle x,e_i\rangle e_i$.  Then
$x-y_F\perp y_F$, so
\[
  \|x\|^2=\|x-y_F\|^2+\sum_{i\in F}|\langle x,e_i\rangle|^2.
\]
Take the supremum over finite $F$.
\end{proof}

\begin{notetheorem}[Closed span]
The closed span of an orthonormal family is
\[
  \overline{\Span\{e_i:i\in I\}}
  =\left\{\sum_{i\in I}a_ie_i:
       \sum_i|a_i|^2<\infty\right\}.
\]
For $x$ in this closed span,
\[
  x=\sum_{i\in I}\langle x,e_i\rangle e_i.
\]
\end{notetheorem}

\begin{notetheorem}[Maximal orthonormal families]
An orthonormal family $(e_i)$ is maximal if and only if its closed linear span
is all of $H$.
\end{notetheorem}

\begin{proof}
If the closed span is proper, choose a nonzero vector in its orthogonal
complement and normalize it; adjoining that vector enlarges the family.
Conversely, a vector orthogonal to a dense span is zero.
\end{proof}

\begin{notedefinition}[Orthonormal basis]
An orthonormal family with dense linear span is an orthonormal basis.
\end{notedefinition}

\begin{notetheorem}[Parseval identities]
If $(e_i)_{i\in I}$ is an orthonormal basis, then for $x,y\in H$,
\[
  x=\sum_{i\in I}\langle x,e_i\rangle e_i,
\]
\[
  \langle x,y\rangle
  =\sum_{i\in I}\langle x,e_i\rangle
      \overline{\langle y,e_i\rangle},
\]
and
\[
  \|x\|^2=\sum_{i\in I}|\langle x,e_i\rangle|^2.
\]
\end{notetheorem}

\begin{notetheorem}[Totality]
A subset $M\subseteq H$ is total if and only if $M^\perp=\{0\}$.  If $H$ is
complete, this is equivalent to $\overline{\Span M}=H$.
\end{notetheorem}

\begin{notetheorem}[Orthonormal families in separable spaces]
Every orthonormal family in a separable Hilbert space is countable.  Every
separable Hilbert space has a countable orthonormal basis.
\end{notetheorem}

\begin{proof}
Distinct orthonormal vectors have distance $\sqrt2$.  Balls of radius
$1/3$ around them are disjoint, and each contains a distinct member of a
fixed countable dense set.  Thus the family is countable.  For existence,
apply Gram--Schmidt to a dense sequence after discarding vectors already in
the span of their predecessors.
\end{proof}

\begin{notetheorem}[Classification by Hilbert dimension]
Two Hilbert spaces are isometrically isomorphic if and only if their
orthonormal bases have the same cardinality.
\end{notetheorem}

\begin{proof}
First, any two orthonormal bases of one Hilbert space have the same
cardinality.  Indeed, if $E$ and $F$ are such bases, then for each $e\in E$
the set $\{f\in F:\langle e,f\rangle\ne0\}$ is countable by Bessel's
inequality.  Totality of $E$ shows that the union of these countable sets is
all of $F$, so $|F|\le |E|$ for infinite bases; reversing the roles gives
equality, while the finite-dimensional case is elementary.

Given equally indexed bases $(e_i)_{i\in I}$ and $(f_i)_{i\in I}$, define
\[
  T\!\left(\sum_ia_ie_i\right)=\sum_ia_if_i.
\]
Parseval shows that $T$ is a surjective linear isometry.  Conversely, a
surjective linear isometry sends an orthonormal basis to an orthonormal basis,
so the cardinalities agree.
\end{proof}

\begin{notetheorem}[Projection and Fourier coefficients]
Let $(e_i)_{i\in I}$ be an orthonormal basis of a closed subspace $M$.  Then
\[
  P_Mx=\sum_{i\in I}\langle x,e_i\rangle e_i.
\]
If $\varphi\in H^*$ and $\varphi(e_i)=a_i$, then
$(a_i)\in\ell^2(I)$ and the Riesz representative is
\[
  h=\sum_{i\in I}\overline{a_i}e_i,
\]
with $\|\varphi\|^2=\sum_i|a_i|^2$ under the linear-first convention.
\end{notetheorem}

\begin{proof}
The difference between $x$ and the displayed sum is orthogonal to every basis
vector and hence to $M$.  The functional statement follows by applying the
Riesz theorem and Parseval to its representing vector.
\end{proof}

\begin{noteexercise}[No nonzero vector orthogonal to a total family]
If $M$ is total and $x\perp M$, then $x=0$.
\end{noteexercise}

\begin{proof}
Orthogonality to $M$ implies orthogonality to its closed span, which is all of
$H$; in particular $\langle x,x\rangle=0$.
\end{proof}
